\section{Motivation}
\label{chap:motivation}

\subsection{Ausgangsproblem}

Ein Muttermal fällt auf, und plötzlich stellt sich die Frage: Sah das schon immer so aus? Ohne ältere Aufnahmen lässt sich das aus der Erinnerung oft schwer beantworten. Einzelne Fotos in der Smartphone-Galerie helfen nur bedingt, wenn sie zwischen anderen Bildern verschwinden oder später nicht mehr klar ist, welche Körperstelle fotografiert wurde. Genau hier setzt unsere Projektidee an: Hautstellen sollen sich einfach dokumentieren und über einen längeren Zeitraum wiederfinden lassen.

Neben der Dokumentation geht es um den Umgang mit Unsicherheit. Wer eine Hautstelle fotografiert, möchte ein verständliches Ergebnis und wissen, was damit gemeint ist. Die geplante KI-Risikoeinstufung soll deshalb nicht allein stehen, sondern durch eine Erklärung ihrer Grenzen und Hinweise zum weiteren Vorgehen ergänzt werden. Im Studienprojekt untersuchen wir, wie sich dieser Ablauf technisch und in der Benutzeroberfläche umsetzen lässt. Eine verlässliche medizinische Beurteilung können wir mit dem Prototyp nicht anbieten.

\subsection{Nutzen der App}

Unser Ziel ist es, die Beobachtung der eigenen Haut übersichtlicher und im Alltag einfacher zu machen. Dafür verbindet die App mehrere Schritte, die sonst getrennt stattfinden würden:

\begin{itemize}
\item \textbf{Hautstellen gezielt festhalten:} Eine geführte Kameraaufnahme unterstützt beim Fotografieren. Die Qualitätsprüfung soll ungeeignete Aufnahmen erkennen und zu einer erneuten Aufnahme auffordern.
\item \textbf{Aufnahmen wiederfinden:} Fotos werden einer Hautstelle zugeordnet und mit ihrem Aufnahmedatum gespeichert. Eine einfache Körperübersicht hilft bei der Orientierung.
\item \textbf{Einen Verlauf aufbauen:} Frühere Aufnahmen bleiben verfügbar. Nutzer müssen sich dadurch nicht allein auf ihre Erinnerung verlassen. Eine automatische Erkennung von Veränderungen ist damit noch nicht gemeint.
\item \textbf{Beobachtung nicht vergessen:} Einstellbare Erinnerungen unterstützen dabei, die Dokumentation fortzuführen.
\end{itemize}

Das Smartphone bietet sich dafür an, weil Kamera, Speicher und Benachrichtigungen bereits vorhanden sind. Zusätzliche Hardware ist für den geplanten Ablauf nicht vorgesehen. Der Verzicht auf ein Nutzerkonto und die zunächst lokale Speicherung sollen den Einstieg erleichtern. Gleichzeitig ist es bei persönlichen Hautfotos wichtig, dass Nutzer nachvollziehen können, wo ihre Bilder gespeichert werden.

\subsection{Motivation im Studienprojekt}

Für das Modul Mobile Software Engineering ist das Thema interessant, weil es mehrere typische Aufgaben mobiler Entwicklung in einem zusammenhängenden Anwendungsfall verbindet: Kamerazugriff, lokale Datenhaltung, Benachrichtigungen und die Einbindung eines KI-Modells auf dem Gerät. Dabei reicht es nicht, dass die einzelnen Funktionen technisch funktionieren. Auch die Bedienung muss verständlich sein, besonders bei Aufnahmefehlern oder einer nicht möglichen Auswertung.

\textbf{Der Schwerpunkt liegt auf einem nachvollziehbaren mobilen Prototyp.} Wir wollen zeigen, wie Aufnahme, experimentelle Auswertung und Dokumentation zusammenspielen. Ob die KI medizinisch zuverlässige Ergebnisse liefert, lässt sich daraus nicht ableiten. Die App ersetzt keine ärztliche Untersuchung; auch eine niedrige angezeigte Risikostufe darf nicht als Entwarnung dargestellt werden.
